\begin{frame}{¿Qué es un ratio financiero?}
\color{fundamentos}
\textbf{Relación entre partidas contables}
\color{black}

\vspace{0.7em}
Un ratio financiero compara dos magnitudes de los estados financieros para evaluar una dimensión concreta: pago, eficiencia, endeudamiento o rentabilidad.

\vspace{0.7em}
\begin{block}{Idea central}
El valor numérico solo adquiere sentido cuando se interpreta junto con el sector, la política de la empresa, periodos anteriores y objetivos de gestión.
\end{block}
\end{frame}

\begin{frame}{Importancia y limitaciones}
\begin{columns}[T]
\column{0.48\textwidth}
\textbf{Importancia}
\begin{itemize}
  \item Facilita comparación.
  \item Detecta tendencias.
  \item Resume información contable.
\end{itemize}
\column{0.48\textwidth}
\textbf{Limitaciones}
\begin{itemize}
  \item Requiere contexto.
  \item No se interpreta aisladamente.
  \item Depende de calidad contable.
\end{itemize}
\end{columns}
\end{frame}

\begin{frame}{Estados financieros utilizados}
\begin{itemize}
  \item \textbf{Estado de Situación Financiera:} activos, pasivos y patrimonio.
  \item \textbf{Estado de Resultados:} ventas, costos, gastos y utilidad.
  \item \textbf{Estado de Flujos de Efectivo:} útil para contrastar liquidez real.
\end{itemize}

\vspace{0.8em}
\begin{center}
\begin{tikzpicture}[node distance=1.1cm, every node/.style={rounded corners, draw, align=center, minimum width=3cm, minimum height=0.8cm}]
\node (esf) {Situación\\Financiera};
\node[right=of esf] (er) {Resultados};
\node[right=of er] (efe) {Flujos de\\Efectivo};
\draw[-{Stealth}] (esf) -- (er);
\draw[-{Stealth}] (er) -- (efe);
\end{tikzpicture}
\end{center}
\end{frame}

\begin{frame}{Clasificación de ratios}
\begin{center}
\begin{tabular}{ll}
\toprule
Categoría & Evalúa \\
\midrule
Liquidez & Capacidad de pago de corto plazo \\
Gestión & Eficiencia del ciclo operativo \\
Solvencia & Participación de deuda y patrimonio \\
Rentabilidad & Generación de beneficios \\
\bottomrule
\end{tabular}
\end{center}
\end{frame}
